%%%%%%%%%%%%%%%%%%%%%%%%%%%%%%%%%%%%%%%%%%%%%
\section{Conclusions}
\label{Sec:Conclusions}
%%%%%%%%%%%%%%%%%%%%%%%%%%%%%%%%%%%%%%%%%%%%%

In these notes we have studied numerically the evolution of small perturbations of the self-consistent steady states of a spherically symmetric collisionless gas which is trapped in an external central potential, in the reduced model in which all the gas particles have the same magnitude of the angular momentum. Our analysis is based on three tools which agree with each other: a PIC code whose particles are placed on a regular lattice in the action-angle variables of the steady state and carry the perturbation in their weights, a solution of the linearized equations in the same variables, and the gain $\lambda(\omega)$ of the self-consistent loop. In this section we summarize the results and list some open problems.

Below the threshold mass at which $\lambda_{edge} = 1$, the perturbation decays. When the self-gravity of the gas is not too weak, the decay is dominated at intermediate times by a damped oscillation, the analogue of Landau damping, whose frequency lies inside the frequency band and approaches its lower edge, while the damping rate decreases, as the threshold is approached. At late times the decay is algebraic. In a fixed potential, a DF which vanishes like $(E_0 - E)^k$ at the edge of its support gives rise to a tail proportional to $t^{-(k+1)}$; the self-gravity of the perturbation reduces this exponent to $-k$ after the time $t_\times$, which we have estimated to leading order in the mass of the gas. For the polytropes with $M_{gas}/M = 0.01$ this time exceeds the time span of our simulations, whereas for $k = 2$ and $M_{gas}/M = 1$ the exponent of the tail is $-2.07$ in the interval $2000\leq t\leq 4000$ and $-2.04$ in $4000\leq t\leq 8000$.

The threshold mass depends strongly on the width of the frequency band and on the behavior of the DF at the edge of its support. For the lowered Maxwellians in the isochrone potential it lies between $0.04$ and $0.27$ times the mass which generates the external potential, depending on the width of the support of the DF, whereas for the polytropes in the Kepler potential it is $0.514$, $0.920$, $1.53$ and $2.38$ for $k = 1.25$, $1.5$, $2$ and $3$, respectively. Hence, the damping of the perturbations which has been established for small masses in Ref.~\cite{mHgRmScS25} does not extend beyond these values. For $k\leq 1$ an oscillating mode exists for any mass, but for small masses its frequency lies so close to the edge of the band that it cannot be observed in a time evolution: for $k = 0.75$ and $M_{gas}/M = 0.01$ the distance from the edge is $2\times 10^{-8}$ times the width of the band. Above the threshold the perturbation contains an oscillating mode, and its mixing is not completed. The frequency of the mode which follows from the equation $\lambda(\omega_d) = 1$ agrees with the one of the linearized solutions to within $2\times 10^{-4}$ in relative terms, and with the one of the PIC simulations to within $5\times 10^{-5}$ for small amplitudes.

The linear regime of a mode which lies close to the edge of the band is restricted to small amplitudes. Above the critical amplitude~(\ref{Eq:epsc}), which vanishes quadratically with the distance of the mode from the band, the gas particles near the edge of the steady state are trapped by the mode. A pendulum model gives the largest action variable which is reached by these particles to within $12$ percent, without adjustable parameters, and the trapping reduces the amplitude of the mode, which then oscillates with a period of about $1.4$ times the bounce period. Below the threshold, the same mechanism stops the damping of a perturbation of large amplitude, as predicted by O'Neil for plasma waves~\cite{tO65}.

The numerical methods required some care. The observables have to be computed with the action-angle variables of the total potential, since those of the external potential leave a time-independent residue which could be mistaken for an undamped mode. The quiet start on a lattice in action-angle variables reduces the noise by several orders of magnitude, but for $M_{gas}/M = 1$ its regularity is lost after $t\approx 1000$ with $10^4$ particles, and about $10^5$ particles and amplitudes $\varepsilon\leq 0.01$ are needed in order to reproduce the solutions of the linearized equations within a few percent.

Several questions remain open. The derivation of the tail $t^{-k}$ is formal, and its amplitude is only known to leading order in $M_{gas}/M$; the Laplace representation of Refs.~\cite{jBaOyY11,jByY13} would also give the amplitude beyond this order. The depth of the loss of amplitude which is caused by the trapped particles is not explained by the pendulum model. The simulations with $M_{gas}/M = 1$ and $k = 0.75$, $1$ and $2$ have to be repeated with $10^5$ particles and smaller amplitudes. Finally, the extension to a distribution of angular momenta, for which the frequency band and its edge depend also on $L$, is left for future work. \EJ{To be completed after the discussion with Olivier; reference to the companion article on initial data far from equilibrium to be added.}
